De profeet zocht een stad.
Wat hij gekozen had,

was eentje met een bron,
die mensen vangen kon.

Waar nevels opstijgen
en zoetjes neerzijgen

saam met zware regen
als een mistig zegen.

Een smaragd-groene kleur
als ook een zwavelgeur

had water in een put
bij een arbeidershut.

Je hoorde er echo’s.
Al waren geesten boos.

Dat was de juiste plek
voor zijn heilige stek.

Hij liet volgelingen
de veengrond doordringen

tot de bron er doorbrak.
Het water smaakte brak

en was aangenaam warm.
De arbeider was arm

en verkocht graag zijn hut.
Niet wetend van zijn nut

als potentiële kerk.
Men daar aan een bouwwerk

met veel passie begon
rondom een warme bron.

Nu niet zo bijzonder.
Toen zag men een wonder.

De kerk werd ermee rein,
als vrijplaats voor welzijn,

aldus de heilsprofeet.
Het heetste water deed

dienst als vloerverwarming.
Toch een verademing

in vergelijking met
het antivriespakket
dat men mee in kerken
bracht om zich te sterken

tegen kou en kilte
in de morgenstilte.

Ook aangesproken door de piëtistische beweging.

Jezus sprak over valse profeten,
die de kudde op het verkeerde been zetten.
Zij verschijnen als wolven in schaapskleren,
maar hun daden zijn slecht en hun woorden verteren.

Toch is er een leider, een sekteleider,
die velen volgen als een echte bevrijder.
Hij spreekt in tongen en doet vele wonderen,
zijn volgelingen zijn vol lof en tonen hun zonden.

Ze worden weer toegelaten in de stad,
waar ze eerder werden verbannen als een misdadigerklad.
Maar Jezus zei: "Bewijs hen aan hun vruchten,
want valse profeten zullen uiteindelijk vluchten."

Dus wees gewaarschuwd voor degenen die komen,
die beweren dat ze de waarheid hebben genomen.
Bekijk hun daden en luister goed naar hun woorden,
want alleen dan zul je weten of ze echt zijn of verworden.

In de beginjaren van de 18e eeuw was er nog geen uitgebreid nationaal politie- of rechtssysteem in Noord-Duitsland. Hierdoor konden criminelen gemakkelijker tussen steden reizen en in andere steden schuilen als ze gezocht werden in een specifieke stad. Er was weinig communicatie en coördinatie tussen de verschillende lokale autoriteiten en dus konden criminelen vrij gemakkelijk ontsnappen aan de wet.

Dat gezegd hebbende, waren er natuurlijk wel risico's verbonden aan het verplaatsen van stad naar stad als crimineel. Er was altijd een kans dat een lokale autoriteit informatie zou delen met andere steden of zelfs met andere landen, waardoor criminelen uiteindelijk alsnog gevangen konden worden genomen. Bovendien zou een crimineel zich nog steeds moeten verbergen voor lokale wetshandhavingsinstanties, waardoor hij of zij zich waarschijnlijk niet volledig vrij zou voelen.

Over het algemeen was het dus mogelijk voor criminelen om in andere steden te schuilen als ze gezocht werden, maar het was niet zonder risico's en er waren zeker uitdagingen en beperkingen aan verbonden.

In de beginjaren van de 18e eeuw waren er geen formele uitleveringsverdragen tussen steden of landen, en daarom had elke stad het recht om te beslissen wie ze als inwoner toelieten. Als een crimineel bijvoorbeeld werd geaccepteerd als bewoner van een andere stad, dan zou het inderdaad moeilijk zijn voor de oorspronkelijke stad om die crimineel te laten uitleveren.

Er was echter wel een zekere mate van samenwerking tussen de lokale autoriteiten, en als een stad wist dat een crimineel zich elders in een andere stad bevond, konden ze samenwerken om de crimineel op te sporen en te arresteren. Dit hing echter sterk af van de lokale politieke en sociale verhoudingen, en sommige steden zouden waarschijnlijk meer bereid zijn om samen te werken dan andere.

Kortom, hoewel een stad het recht had om een crimineel als inwoner toe te laten en het in sommige gevallen moeilijk zou zijn om de crimineel uit te leveren, was er nog steeds een kans dat de oorspronkelijke stad zou samenwerken met andere autoriteiten om de crimineel op te sporen en te arresteren.

Zonder domheid is er geen leren mogelijk. Leren => wijsheid => macht.
De sekteleider echter kon het niet hebben die domheid en zeker niet als anderen hem een spiegel van zijn eigen onwetendheid voorhouden. Juist als hij wist dat het waar was, ging hij die ander die vaak het goedbedoelde aanvallen.

Toeval het valt je toe => alles onder controle

De valse profeet

Hier het verhaal hoe de stad is overgenomen door een religieuze sekte die gevluchte criminelen of vogelvrijen of andere kwetsbaren heeft geworven.

Domheid of onwetendheid is de basis voor lijden. De positieve variant is dat leren lijden zal verminderen. Hekel aan lemmingen.

Begin 18e eeuw was de macht van de geestelijken in Noord-Duitsland afhankelijk van de specifieke stad en religie. In de lutherse steden was de dominee, ook wel bekend als de predikant, de meest invloedrijke geestelijke. Hij was verantwoordelijk voor de leiding van de kerk en voor het geven van preken. De dominee was meestal een zeer gerespecteerde figuur in de gemeenschap en had vaak veel invloed op de lokale politiek.

In katholieke steden was de bisschop of de pastoor de belangrijkste geestelijke autoriteit. De bisschop was verantwoordelijk voor het bestuur van de bisdommen en de pastoor leidde de lokale parochie. Beide figuren hadden vaak veel invloed op de lokale politiek en de sociale omgang van de gemeenschap.

De macht van deze geestelijken varieerde afhankelijk van verschillende factoren, waaronder hun persoonlijke populariteit, rijkdom en connecties. In sommige steden konden geestelijken een aanzienlijke politieke macht uitoefenen, terwijl in andere steden hun invloed minder groot was.

Het is belangrijk op te merken dat in de 18e eeuw de macht van de geestelijkheid afnam, met name in protestantse steden, als gevolg van de Verlichting en de secularisatiebeweging. Toch hadden geestelijken in die tijd nog steeds aanzienlijke invloed op de gemeenschappen waarin ze werkten.

Ja, er waren inderdaad bisschoppen en pastoren in de 18e eeuw die heimelijk affiniteit hadden met lutherse ideeën, maar die dit vanwege hun rol niet openlijk naar buiten brachten of zich plotseling bekeerden. Dit kwam met name voor in gebieden waar er een religieuze strijd was tussen verschillende confessionele groepen, zoals tussen de katholieken en protestanten.

Sommige bisschoppen en pastoren hadden persoonlijke twijfels over de katholieke leer en praktijken en voelden zich meer aangetrokken tot de leer van Luther en andere protestantse reformators. Sommigen waren echter bang om openlijk hun geloof te veranderen vanwege de mogelijke gevolgen voor hun positie of zelfs hun leven.

In plaats daarvan kozen ze ervoor om hun overtuigingen in het geheim te praktiseren of om hun leerlingen en volgelingen subtiel te beïnvloeden in de richting van lutherse ideeën. Soms leidde dit tot onrust binnen de kerk of gemeenschap, vooral als deze leiders hun overtuigingen toch uiteindelijk naar buiten brachten of bekeerden tot het lutheranisme.

Een voorbeeld van een dergelijke leider is Nikolaus Ludwig von Zinzendorf, een graaf en bisschop uit de 18e eeuw die een belangrijke rol speelde in de Moravische kerk, een protestantse denominatie die ontstond uit de hervormingsbeweging in Bohemen en Moravië. Zinzendorf had sympathieën voor lutherse ideeën en werd beschuldigd van het bevorderen van ketters denken en praktijken binnen de Moravische kerk. Hij moest daarom aftreden als bisschop, maar hij bleef actief in de kerk en bleef lutherse ideeën verspreiden.

Er was eens een geestelijke in de stad
Die steeds meer van een sekte hield
Een leider die vogelvrijen had gebracht
En beschuldigde heksen, met weinig geld

Hij sprak over God die van allen houdt
En dat hun zonden zouden worden vergeven
De geestelijke vond dat dit waarheid was
En dat de sekte mensen een nieuw leven kon geven

De sekte leider was charmant en slim
Hij overtuigde zijn volgelingen met gemak
En de geestelijke zag zijn idealen in hem
En was bereid om de sekte te steunen met zijn kracht

Maar de sekte leider had zijn eigen plan
Hij gebruikte zijn volgelingen voor zijn doel
En de geestelijke zag zijn fouten niet, als fan
Totdat de sekte leider zijn eigen graf groef in zijn rol

De geestelijke zag toen de waarheid van de sekte
Dat het meer was dan alleen maar liefde en genade
En hij moest toegeven dat hij een fout had gemaakt
En beloofde dat hij nooit meer een sekte zou omarmen met waarden als daden.

=> Eersy werd de sekte hard onderdrukt door de katholieke kerk en werdt de sekteleider opd de brandstapel gezet. Iets wat hij had aangekondigd, maar eigenlijk ook bewust opriep om zijn boodschap krachtiger te maken. Dat leidde alleen maar publieke verontwaardiging en meer steun. Ook van de bisschop die steunde de overname van de stad door de sekte en berichtte zijn superieuren dat juist de sekte de ware katholieken waren. Aangevuld met een flinke som geld. De stedelingen die niet waren aangesloten werden verbannen of weggepest. Sommigen leden zagen dit en andere praktijken ingaan tegen de leer (God houdt van iedereen) en begonnen te rebelleren. Dat leidde tot harde vervolging door sekteleden. Dit is te vergelijken met

In het Noorden van Duitsland,
Daar ontstond een sekte, geen band,
Voor hen die uitgesloten waren,
Criminelen, heksen, zonder baren.

Een charismatische leider kwam,
Hij sprak over een goddelijk plan,
Dat God alle mensen liefheeft,
En dat zijn boodschap de juiste is, die geneest.

Hij dichtte in gekruist rijm,
Om zijn boodschap te verspreiden, fijn,
Zijn woorden raakten de harten,
En zijn volgelingen bleven niet te sparten.

Een sekteleider weet vaak goed,
Hoe hij mensen aan zich bindt, zo zoet,
Hij speelt in op hun diepste verlangens,
En belooft hen de hemel, als in vangens.

Hij biedt hen een oplossing,
Voor hun problemen, zo'n inlossing,
Een nieuwe familie, een gemeenschap,
Waar men zich thuis voelt, zo fijn, zo zacht.

De leider manipuleert en indoctrineert,
Hersenspoelt en isoleert, zo leert,
Hij maakt gebruik van psychologische trucs,
En wakkert angsten aan, met een zucht.

Maar waarom lukt het hem om dit te doen?
Waarom geloven mensen, wat is hun zoen?
Vaak is het zo dat mensen zoeken,
Naar zingeving en verbondenheid, om te vloeken.

En als de sekte de enige plek lijkt,
Waar zij dit vinden, hoe uniek,
Dan zijn ze bereid alles te doen,
Om bij de sekte te horen, het te zoen.

De sekteleider, zo blijkt keer op keer,
Speelt in op deze behoeften, het is geen keer,
Hij weet wat mensen zoeken, en wat hen drijft,
En zorgt dat de sekte in hun hart beklijft.

De sekte groeide snel en sterk,
De leider kreeg het respect, de eer, de werk,
Maar achter de glans en pracht,
Gingen geheimen schuil, zo zacht.

De sekteleider had beloofd,
Een hernieuwde erkenning, zo zoet,
Maar de gemeenschap bleef afzijdig,
En erkenning bleef uit, zo bitterlijk.

De leden, teleurgesteld en boos,
Besloten om in opstand te komen, zo broos,
Ze grepen de stad met geweld,
Een daad die tot op heden telt.

Godsdienstsocioloog John G. Gager,
Paste de theorie van Festinger toe,
Op de kwestie Jezus, zo waar,
Een voorspelling die niet uitkwam, het was zo.

De kern van de theorie luidt,
Dat het geloof in de voorspelling juist groeit,
Wanneer het niet uitkomt, hoe raar,
Gelovigen gaan tot extremen, echt waar.

Zo ook bij de sekteleider en zijn volk,
Die de stad met geweld overnamen, zo dol,
En zichzelf overtuigden van hun gelijk,
Ondanks de gevolgen, zo zwaar en rijk.

De sekte viel niet uiteen, zoals verwacht,
Maar nam de stad over, met geweld in hun macht,
Een tragische afloop, zo is het gebleken,
De sekte bleef tot het uiterste strijden, tot het verbleken.
